\documentclass[10pt]{article}
\usepackage[T1]{fontenc}
\usepackage[utf8]{inputenc}
\usepackage[french]{babel}
\usepackage{amsmath,amssymb,mathrsfs,mathpazo}
\usepackage{geometry}
\geometry{a4paper,margin=2cm}
\usepackage{graphicx,array,tabularx,tabularray,colortbl,multirow,multicol,booktabs,enumitem,pifont,ulem,fancybox,diagbox,cancel}
\usepackage{pstricks,pstricks-add,pst-plot,pst-tree,pst-node,pst-text,pst-func,pst-math,pst-bspline,pst-3dplot,pst-eucl,pst-solides3d}
\let\pscalculate\undefined
\usepackage{pst-bezier}
\usepackage{tikz,pgfplots,tkz-tab}
\usetikzlibrary{arrows,arrows.meta,patterns,calc,positioning,trees,decorations.markings,shapes.geometric,shapes.misc,shapes.arrows,intersections,decorations.pathmorphing,decorations.pathreplacing,backgrounds,shadows,babel}
\pgfplotsset{compat=1.18}
\usepgfplotslibrary{fillbetween}
\usepackage[np]{numprint}
\usepackage{listings,makecell,pas-tableur,twoopt}
\let\barre\undefined
\usepackage[pstricks]{tablvar}
\definecolor{darkgreen}{RGB}{0,100,0}
\providecommand{\red}{\color{red}}
\providecommand{\blue}{\color{blue}}
\providecommand{\green}{\color{green}}
\providecommand{\black}{\color{black}}
\providecommand{\white}{\color{white}}
\providecommand{\gray}{\color{gray}}
\providecommand{\cyan}{\color{cyan}}
\providecommand{\magenta}{\color{magenta}}
\providecommand{\yellow}{\color{yellow}}
\providecommand{\violet}{\color{violet}}
\providecommand{\orange}{\color{orange}}
\providecommand{\euro}{\texteuro}
\providecommand{\EUR}{\texteuro}
\providecommand{\ieme}{\textsuperscript{e}}
\providecommand{\degres}{\ensuremath{^\circ}}
\providecommand{\degree}{\ensuremath{^\circ}}
\providecommand{\vv}[1]{\overrightarrow{#1}}
\providecommand{\Coord}[1]{#1}
\providecommand{\decosix}{}
\providecommand{\starredbullet}{\ensuremath{\bullet}}
\pagestyle{empty}
\setlength{\parindent}{0pt}
\usepackage{adjustbox}\DeclareUnicodeCharacter{2236}{:}\DeclareUnicodeCharacter{2190}{\ensuremath{\leftarrow}}
\providecommand{\R}{\mathbb{R}}
\providecommand{\N}{\mathbb{N}}
\providecommand{\Z}{\mathbb{Z}}
\providecommand{\Q}{\mathbb{Q}}
\providecommand{\D}{\mathbb{D}}
\providecommand{\C}{\mathbb{C}}
\providecommand{\e}{{\mathrm{e}}}
\providecommand{\vect}[1]{\overrightarrow{#1}}
\providecommand{\vectt}[1]{\overrightarrow{\mathrm{#1}}}
\providecommand{\vvt}[1]{\overrightarrow{\mathrm{#1}}}
\providecommand{\np}[2][]{{#2~#1}}
\providecommand{\vv}[1]{{\overrightarrow{#1}}}
\renewcommand{\d}{{\mathrm{d}}}
\providecommand{\degres}{{^\circ}}
\providecommand{\degre}{{^\circ}}
\providecommand{\blue}{\color{blue}}
\providecommand{\red}{\color{red}}
\providecommand{\enskip}{\quad}
\providecommand{\llbracket}{\lbrack\!\lbrack}
\providecommand{\rrbracket}{\rbrack\!\rbrack}
\providecommand{\euro}{{\text{€}}}
\providecommand{\no}{{\text{n°}}}
\providecommand{\interval}[2]{{[#1;#2]}}
\providecommand{\intervFO}[2]{{[#1;#2[}}
\providecommand{\dotfill}{\ldots}

\providecommand{\ds}{\displaystyle}
\providecommand{\textding}[1]{#1}
\providecommand{\Oij}{\(\mathrm{O};\vec{\imath},\vec{\jmath}\)}
\providecommand{\Oijk}{\(\mathrm{O};\vec{\imath},\vec{\jmath},\vec{k}\)}
\providecommand{\Ouv}{\(\mathrm{O};\vec{u},\vec{v}\)}
\providecommand{\intervFF}[2]{\ensuremath{[#1;#2]}}
\renewcommand{\euro}{\ifmmode\text{\texteuro}\else\texteuro\fi}
\DeclareUnicodeCharacter{00B0}{\ifmmode^{\circ}\else\textdegree\fi}
\DeclareUnicodeCharacter{20AC}{\ifmmode\text{\texteuro}\else\texteuro\fi}
% Énoncés : Éducation nationale. Transcription et corrigés : APMEP, Denis Vergès et contributeurs. Découpage et classement : CoopMaths, Sébastien Lozano, Éric Elter et contributeurs. Adaptation de la présentation : Maths Explorer.
% CC BY-SA : https://coopmaths.fr/www/about/
% https://coopmaths.fr/alea/static/bac/2025/tex/bac_2025_05_sujet2_ameriquenord_secours_3_cor.tex
\begin{document}
\noindent\textit{APMEP / CoopMaths. CC BY-SA. Présentation adaptée par Maths Explorer.}\par\medskip
\medskip

\begin{enumerate}
\item Ici, on interroge une personne au hasard, donc toutes les personnes qui constituent la population du pays ont la même probabilité d'être choisis : c'est une situation d'équiprobabilité et donc les proportions sont assimilables à des probabilités.

On a donc :
\begin{itemize}
\item $P(C) = 0,02$, car 2~\% de la population du pays a été contaminée;
\item $P(V) = 0,9$, car 90~\% de la population a été vaccinée;
\item $P_C(V) = 0,62$, car 62~\% des personnes contaminées ont été vaccinées.
\end{itemize}

\item 
	\begin{enumerate}
		\item On a :\quad $P(C \cap V) = P(C) \times P_C(V) = 0,02 \times 0,62 = 0,0124$.
		\item Les évènements $C$ et $\overline{C}$ partitionnent l'univers, donc, d'après la formule des probabilités totales :

$P(V) = P(C \cap V) + P\big(\,\overline{C}\cap V\big) \iff P\big(\,\overline{C} \cap V\big) = P(V) -
 P(C \cap V)$

On a donc : \quad $P\big(\,\overline{C} \cap V\big) = P(V) - P(C \cap V) = 0,9 - 0,0124 = \np{0,8876}$
	\end{enumerate}
\item Pour compléter l'arbre de probabilité, il nous faut $P_{\overline{C}}(V)$.

Par définition :\quad $P_{\overline{C}}(V) = \dfrac{P\big(\,\overline{C} \cap V\big)}{P\big(\,\overline{C}\,\big)} = \dfrac{0,8876}{1 - P(C)}=\dfrac{0,8876}{0,98} \approx 0,9057$ au dix-millième près.

Et donc : \hfill~
\begin{tikzpicture}[xscale=1,yscale=1,baseline={(R.base)}]

\tikzstyle{fleche}=[->,>=latex,thick]
\tikzstyle{noeud}=[fill=white,circle,inner sep=2pt]
\tikzstyle{feuille}=[fill=white,circle,inner sep=2pt]
\tikzstyle{etiquette}=[pos=0.6,fill=white, inner xsep=3pt, inner ysep=1.5pt]

\def\DistanceInterNiveaux{4.5}
\def\DistanceInterFeuilles{0.9}

\def\NiveauA{(0)*\DistanceInterNiveaux}
\def\NiveauB{(1)*\DistanceInterNiveaux}
\def\NiveauC{(2)*\DistanceInterNiveaux}
\def\InterFeuilles{(-1)*\DistanceInterFeuilles}

\node[noeud] (R) at ({\NiveauA},{(1.5)*\InterFeuilles}) {};
\node[noeud] (Ra) at ({\NiveauB},{(0.5)*\InterFeuilles}) {$C$};
\node[feuille] (Raa) at ({\NiveauC},{(0)*\InterFeuilles}) {$V$};
\node[feuille] (Rab) at ({\NiveauC},{(1)*\InterFeuilles}) {$\overline{V}$};
\node[noeud] (Rb) at ({\NiveauB},{(2.5)*\InterFeuilles}) {$\overline{C}$};
\node[feuille] (Rba) at ({\NiveauC},{(2)*\InterFeuilles}) {$V$};
\node[feuille] (Rbb) at ({\NiveauC},{(3)*\InterFeuilles}) {$\overline{V}$};

\draw[fleche] (R.east)--(Ra.west) node[etiquette] {$0,02$};
\draw[fleche] (Ra.east)--(Raa.west) node[etiquette] {$0,62$};
\draw[fleche] (Ra.east)--(Rab.west) node[etiquette] {$0,38$};
\draw[fleche] (R.east)--(Rb.west) node[etiquette] {$0,98$};
\draw[fleche] (Rb.east)--(Rba.west) node[etiquette] {$0,9057$};
\draw[fleche] (Rb.east)--(Rbb.west) node[etiquette] {$0,0943$};
\end{tikzpicture}
\hfill~

\item Par définition :\quad $P_V(C) = \dfrac{P(V \cap C)}{P(V)} = \dfrac{0,0124}{0,9} \approx 0,0138$.

Dans le contexte de l'exercice, cela signifie qu'environ 1,38~\% des personnes vaccinées ont été contaminées par le virus.

\item
	\begin{enumerate}
		\item \og Parmi les personnes non contaminées, il y a dix fois plus de personnes vaccinées que de personnes non vaccinées.\fg{}.

Cette affirmation est \textbf{fausse}, car un peu exagérée, parmi les personnes non contaminées, environ 90,57~\% d'entre elles sont vaccinées, quand 9,43~\% d'entre elles ne le sont pas, mais le décuple de 9,43~\% est 94,3~\%, qui est donc supérieur à 90,57~\%.

Il y a $\dfrac{90,57}{9,43}\approx 9,6$ fois plus de personnes vaccinées que non vaccinées parmi les personnes non contaminées.

		\item \og Plus de 98~\% de la population vaccinée n'a pas été contaminée.\fg{}

Cette affirmation est \textbf{vraie}, car on a calculé qu'environ 1,38~\% des personnes vaccinées ont été contaminées par le virus, donc, par complémentaire, cela implique qu'environ 98,62~\% des personnes vaccinées n'ont pas été contaminées. Il est donc correct de dire que c'est plus de 98~\%.
	\end{enumerate}
\item
	\begin{enumerate}
		\item 
\begin{itemize}
\item On a une épreuve aléatoire à deux issues. On qualifie de succès l'évènement $C$, de probabilité $p = 0,02$;
\item Cette épreuve est répétée $n =20$ fois pour constituer l'échantillon de 20 personnes. La répétition étant assimilable à des tirages successifs avec remise, on peut dire que les répétitions sont identiques et indépendantes;
\item La variable aléatoire $X$ compte le nombres de personnes contaminées, donc le nombre de succès, dans l'échantillon.
\end{itemize}

Ces trois éléments garantissent que $X$ suit la loi binomiale, de paramètres $n = 20$ et $p = 0,02$.

		\item La probabilité que 4 personnes exactement soient contaminées dans ce groupe de 20 personnes est :

			$P(X = 4) = \displaystyle \binom{20}{4}\times 0,02^4 \times 0,98^{20 - 4} = \np{4845}\times 0,02^4 \times 0,98^{16} \approx 5,6\times 10^{-4}$

Finalement, avec les consignes d'arrondi : \quad $P(X = 4) \approx 0,0006$
	\end{enumerate}
\end{enumerate}

\medskip

\end{document}
